%!TeX encoding = UTF-8 Unicode
%!TeX spellcheck = en-UK
% !TeX root = Perrinet26mnesis.tex
% -*- coding: UTF-8; -*-
\newcommand{\Title}{Working Memory with Polychronous Chains}
\newcommand{\ShortTitle}{MNESIS: Memory with Polychronous Chains}% running head
\newcommand{\Keywords}{spiking neural networks; polychronisation; heterogeneous delays; working memory.}
\newcommand{\FirstLP}{Laurent U}
\newcommand{\LastLP}{Perrinet}
\newcommand{\AuthorLP}{\FirstLP\ \LastLP}
\newcommand{\EmailLP}{laurent.perrinet@cnrs.fr} %univ-amu.fr}
\newcommand{\orcidLP}{0000-0002-9536-010X}
\newcommand{\Department}{Institut de Neurosciences de la Timone (UMR 7289)}%
\newcommand{\Affiliation}{Aix Marseille Univ, CNRS}%
\newcommand{\Street}{27 boulevard Jean Moulin}%
\newcommand{\PostCode}{13005}%
\newcommand{\City}{Marseille}%
\newcommand{\Country}{France}%
\newcommand{\WebsiteLP}{https://laurentperrinet.github.io}%
\newcommand{\Acknowledgements}{
This work was performed using HPC resources from GENCI-IDRIS (Grant 2025--AD010314955R2).
The authors declare no competing interests regarding this manuscript. The funders had no role in the design of the study; in the collection, analyses, or interpretation of data; in the writing of the manuscript; or in the decision to publish the results.
For the purpose of open access, the authors have applied a CC-BY public copyright licence to any Author-Accepted Manuscript version arising from this submission, and the different versions are available on arXiv at \url{https://arxiv.org/abs/2604.14096}.
}
\newcommand{\AIstatement}{
During the preparation of this paper, the Gemma4:31b generative AI tool (\url{https://ollama.com/library/gemma4:31b-mlx}) and Claude Opus 5.5 were used solely for language editing---improving the clarity, grammar, and phrasing of drafted text---and for limited technical assistance with code development. These tools were not used for generating research ideas, formulating research questions or hypotheses, designing the methodology, conducting or interpreting data analysis, or performing mathematical derivations or calculations. All intellectual contributions, analytical work, and mathematical reasoning in this paper are the authors' own. All AI-assisted content has been reviewed by the authors, who assume full responsibility for the accuracy and integrity of the entire work.
}
\newcommand{\Abstract}{%
Working memory --- the ability to store and recall precise temporal patterns of activity --- remains an open challenge for spiking neural networks (SNNs). A recurrent SNN of $N=1024$ neurons is proposed, in which each synapse is equipped with $D = 41$ heterogeneous delays, parameterised as a weight tensor $\mathbf{W} \in \mathbb{R}^{N \times N \times D}$ and trained end-to-end with surrogate-gradient backpropagation through time. Each stored pattern is represented as a sequential chain of overlapping Spiking Motifs: contiguous context windows of length $D$ that uniquely predict the activity at the next time step. A closed-form Hebbian initialisation, derived by deconvolving the LIF membrane response and targeting a membrane potential value above the firing threshold $\vartheta$, already achieves high $F_1$ scores before any gradient step on a benchmark of $M=16$ patterns of duration $T=1000$ steps. With learning, the network reaches a mean $F_1$ score of approximately $0.9$, and retrieval degrades gracefully when the trigger window is corrupted, with high $F_1$ scores maintained up to about $20\%$ bit-flip noise. The network shows attractor-like retrieval dynamics consistent with hippocampal pattern completion and with low-dimensional collective dynamics of population activity on manifolds. These results demonstrate that heterogeneous synaptic delays are an efficient and scalable substrate for working memory in SNNs, with direct implications for neuromorphic edge deployment. %
}
\newcommand{\GeneralPublic}{%
Modelling working memory ---the ability to store and recall precise temporal patterns--- remains an open challenge in computational neuroscience. While recent advances evolved from Long Short-Term Memory (LSTM) networks to transformer architectures, which excel at linking long-range associations, these models deviate from biological plausibility: they are massive and energy-intensive. We present recent work on recurrent spiking neural networks (SNNs), inspired by the structure of information processing in biological systems. Based on the dynamics of realistic yet simple single-neuron models, the approach sheds light on fundamental neural phenomena, such as pattern formation in neural populations and the organisation of high-dimensional brain activity. This work aims to bridge the gap between biologically grounded models and scalable computational solutions, paving the way for more efficient and interpretable AI systems. %
}